% chapter/chapter1.tex
% Mortal Adversaries – Those Who Owe Their Debts in Coin, Steel, or Silence
% Rewritten in the voice of Lerris Fenwood, with marginalia from the Gray Wanderer,
% Saikou Ira, and Dusana of the Raven Road.

\epigraph{“I have seen a man forget his own name because he refused to count the ninth breath. The witness does not forgive; it only remembers.”}{—Lerris Fenwood, marginal note}

The dead are not the only hungry things. I learned this not in the little room with the salt‑ringed table, but in the bright stalls of Thepyrgos, where a clerk sold me a ledger that had been kept by a man who later hanged himself. The sums were perfect. The margin had a single, empty line.

Mortals walk the roads with ledgers in their sleeves and knives in their pockets. They owe debts to patrons, to factions, to the Hollow itself – and they will collect from anyone who crosses their path. This chapter catalogs the living: mercenaries, cultists, faction agents, and those who have simply learned that mercy is a luxury they cannot afford. I have met some of them. I have outlived others.

\section*{Generic Templates}
\epigraph{“They are not evil. They are hungry. There is a difference, though the difference will kill you all the same.”}{—The Gray Wanderer, from a letter I was not meant to read}

The following templates can be scaled using Threat Level (see Appendix). Use them for filler encounters, random patrols, or the quiet before the storm. I have kept the descriptions plain, because terror is best served without ornament.

\subsection*{Bandit Rabble (Threat Level 1)}

\noindent\textbf{Description}: Desperate locals with rusty blades. They flee if half their number falls. Their hunger is for bread, not blood.

\noindent\textbf{Stats}: Body 2, Wits 1, Spirit 1, Presence 1; Melee 1, Athletics 1.\\
\textbf{Harm}: 2. \textbf{Fatigue}: Body 2.

\noindent\textbf{Traits}: \emph{Rabble} – On a Miss when they attack, the GM gains 1 SB (they trip, argue, or hesitate). The sound of their quarrel is often louder than their steel.

\noindent\textbf{Clock}: \emph{Courage [4]} – Ticks each time one of them is wounded. When full, they rout. I have seen it fill on the second death, and on the fourth; there is no accounting for the breaking point.

\gray{A bandit with a full belly is a bad merchant. A bandit with an empty belly is a bad gamble.}

\subsection*{City Watch Recruit (Threat Level 2)}

\noindent\textbf{Description}: Poorly trained but eager. Will call for backup if overwhelmed. Their badge is new; their confidence is older than their sense.

\noindent\textbf{Stats}: Body 2, Wits 1, Spirit 1, Presence 2; Melee 1, Command 1.\\
\textbf{Harm}: 3. \textbf{Fatigue}: Body 2.

\noindent\textbf{Trait}: \emph{Alarm} – On a Partial or Miss, advance the \emph{Reinforcements} timer by 1. The whistle they carry has a note that carries three streets.

\noindent\textbf{Clock}: \emph{Reinforcements [6]} – When full, a sergeant and 1d4 recruits arrive. The sergeant’s face will be the one you remember; the recruits will blur.

\saikou{I have interrogated recruits who did not know they were chasing a ghost. The whistle’s note is the same whether the quarry is flesh or memory.}

\subsection*{Cult Novice (Threat Level 1)}

\noindent\textbf{Description}: True believer, poorly trained. Will gladly die for the cause. Their eyes have the shine of a candle that has been lit from both ends.

\noindent\textbf{Stats}: Body 1, Wits 1, Spirit 2, Presence 1; Lore 1 (cult‑specific), Sway 1.\\
\textbf{Harm}: 2. \textbf{Fatigue}: Body 1.

\noindent\textbf{Trait}: \emph{Faith’s Crumble} – If a player successfully discredits the cult’s doctrine (Presence + Lore vs DV 4), the novice surrenders or flees. The doctrine is often a single sentence; break the sentence, break the believer.

\noindent\textbf{Clock}: \emph{Devotion [4]} – Ticks each time the novice is wounded. When full, they trigger a suicide attack (Harm 2 to everyone in Close). I have seen it fill. I have not forgotten the sound.

\dusana{They are not monsters. They are mothers and fathers who learned the wrong prayer. The prayer is the trap; the believer is the bait.}

\subsection*{Assassin (Threat Level 3)}

\noindent\textbf{Description}: Professional killer. Prefers poison, stealth, and a clean exit. The Veiled Sisters train them. The Black Hand deploys them. You will not see them coming.

\noindent\textbf{Stats}: Body 3, Wits 3, Spirit 2, Presence 2; Stealth 3, Melee 2, Subterfuge 2.\\
\textbf{Harm}: 4. \textbf{Fatigue}: Body 3.

\noindent\textbf{Traits}: \emph{First Strike} – On the first round of combat, the assassin has Dominant Position. They have already counted your breaths.

\noindent\textbf{Traits}: \emph{Poisoned Blade} – On a Critical Hit, the target gains the Condition \emph{Poisoned} (−1 die to physical actions until treated). The poison is not always lethal; sometimes it is only a lesson.

\noindent\textbf{Clock}: \emph{Exit Strategy [4]} – Ticks each round. When full, the assassin disengages and vanishes (cannot be pursued without a relevant Asset). I have watched a door close behind one; when I opened it, the room was empty but for a single salt crystal on the sill.

\lerris{The Black Hand does not recruit from the desperate. They recruit from the precise. An assassin’s ledger is kept in wounds, not ink.}

---

\section*{Faction Operatives}

\epigraph{“The law here flows like water over stone – it finds its own path, and woe to those who try to dam it with paper.”}{—Stone Prince of Aelinnel, quoted by Aqyl, son of Aqyl}

These adversaries are tied to specific factions. Use them when the party has made enemies of the powerful – or when the powerful have taken an interest. I have listed them in order of the danger they present to the soul, not to the body.

\subsection*{Chain‑Lantern Inquisitor (Thepyrgos)}

\noindent\textbf{Description}: Witch‑hunter of the Thepyrgos Synod. Iron mesh hood, brass lantern. The lantern is not for light; it is for testimony.

\noindent\textbf{Stats}: Body 3, Wits 3, Spirit 3, Presence 3; Lore 3, Melee 2, Investigation 2.\\
\textbf{Harm}: 4. \textbf{Fatigue}: Body 3.

\noindent\textbf{Traits}: \emph{Lantern Test} – Once per scene, the Inquisitor may force a PC to test Spirit + Resolve (DV 4) or reveal a surface thought (the GM learns one fact about the PC’s current intent). The thought is never the one the PC would have chosen to share.

\noindent\textbf{Traits}: \emph{Proscription} – If the Inquisitor knows a PC’s true name, they gain +1 die to all social rolls against that PC. Names are keys. The Inquisitor collects keys.

\noindent\textbf{Clock}: \emph{Arrest Warrant [6]} – Ticks each time the party obstructs or lies to the Inquisitor. When full, the Inquisitor returns with 1d4 guards and a sealed writ of seizure. The writ is always stamped with a bell that has no clapper.

\saikou{They burned a woman in Ecktoria for owning a red thread. The thread was for her daughter’s fever. The fire did not discriminate. The Inquisitor filed the report in triplicate.}

\subsection*{Daughter of the Unbroken Cord}

\noindent\textbf{Description}: Mendicant witch, hospice nurse, debt‑collector in gray robes. Accompanied by a Hollowed (a heartless child with silver‑wire eyes). The child does not blink. The child does not age. The child waits.

\noindent\textbf{Stats}: Body 2, Wits 3, Spirit 4, Presence 3; Medicine 2, Lore 2, Sway 2.\\
\textbf{Harm}: 3. \textbf{Fatigue}: Body 2.

\noindent\textbf{Traits}: \emph{Hollowed Witness} – While the Hollowed is within Near, the Daughter counts as having Soft Witness (improve Position by one step for any Rite of the High Mother). The milk‑white eyes record everything; they forgive nothing.

\noindent\textbf{Traits}: \emph{Unwoven Mercy} – When the Daughter would take Harm, she may instead mark a segment on a \emph{Debt Timer } owned by the target (if any). If no debt exists, she may offer a bargain: “I will spare you, but you owe me a name.” The name is never yours to give.

\noindent\textbf{Clock}: \emph{Cord’s Patience [6]} – Ticks each time the Daughter is defied. When full, the High Mother demands payment: a memory, a confession, or a life. I have seen a man pay with the memory of his mother’s face. He did not weep; he could not remember why he should.

\dusana{The soup is free. The price never is. I have eaten their bread. I have not slept well since.}

\subsection*{Ykrul Road‑Judge}

\noindent\textbf{Description}: Tumen scribe with seals from three empires. Settles disputes on the steppe. They carry a bell that never rings; if you hear it, run.

\noindent\textbf{Stats}: Body 3, Wits 4, Spirit 3, Presence 2; Lore (Kon’reh) 3, Command 2.\\
\textbf{Harm}: 4. \textbf{Fatigue}: Body 3.

\noindent\textbf{Traits}: \emph{Kon’reh Binding} – Once per scene, the Road‑Judge may draw a board in the dirt. The party must win a contested Wits + Lore roll (DV 4) or accept a Geas (minor task, no direct harm). The board is always a hexagon; the lines are always straight; the trap is always in the geometry.

\noindent\textbf{Traits}: \emph{Hostage String} – If the Road‑Judge has a token from a PC (hair, blood, a signed name), they gain +1 die to all social rolls against that PC. The token is kept in a leather pouch that smells of iron and old grass.

\noindent\textbf{Clock}: \emph{Foedus Recall [4]} – Ticks each time the party breaks steppe law. When full, the Judge invokes an ancient treaty; the party is declared oath‑breakers and cannot trade at Midh Ahkaz for one season. The grass will not remember your face; the horses will.

\gray{The Kon’reh is not a game. It is a prayer written in angles. Lose the geometry, lose the argument.}

\subsection*{Aeler Breath‑Warden}

\noindent\textbf{Description}: Dwarven engineer who monitors air quality, seals, and keystones. Carries a brass bell and a ledger of breaths. They count to eight; the ninth breath is the mountain’s.

\noindent\textbf{Stats}: Body 4, Wits 3, Spirit 2, Presence 1; Mechanics 3, Craft 2.\\
\textbf{Harm}: 5. \textbf{Fatigue}: Body 4.

\noindent\textbf{Traits}: \emph{Bad Air Pocket} – Once per scene, the Warden may declare a zone “thin air”. Any character in that zone must test Body + Endurance (DV 3) or mark 1 Fatigue. The air tastes of copper and old candle smoke.

\noindent\textbf{Traits}: \emph{Keystone Clause} – If the party is in an Aeler vault, the Warden may seal a door or trigger a collapse by spending 2 SB (advance the \emph{Cave‑In} timer by 2). The clause is always written in the margin; you never see it until the stone groans.

\noindent\textbf{Clock}: \emph{Cave‑In [8]} – When full, the tunnel collapses. All characters in the zone must test Body + Athletics (DV 4) or take Harm 2. The dust settles slowly. The silence after is a judgment.

\lerris{I once watched a Breath‑Warden refuse to open a seal because the requester had miscounted his breaths by one. The requester suffocated. The Warden noted the time of death in his ledger and went to tea.}

\subsection*{Black Banner Mercenary Captain}

\noindent\textbf{Description}: Veteran commander of a free company. Scarred, pragmatic, and always watching the horizon. Loyalty is a coin that fluctuates with the market.

\noindent\textbf{Stats}: Body 4, Wits 3, Spirit 3, Presence 3; Melee 3, Command 3.\\
\textbf{Harm}: 6. \textbf{Fatigue}: Body 4.

\noindent\textbf{Traits}: \emph{Condotta Flip} – Once per scene, the Captain may switch sides. Roll a die: on 1–3, the Captain and 1d4 mercenaries betray the party; on 4–6, they demand double pay. The negotiation is always polite, even when the knives are drawn.

\noindent\textbf{Traits}: \emph{Payday} – At the start of each scene involving the Captain, the GM may tick the \emph{Arrears} timer. When it fills, the Captain’s company mutinies (treat as 1d6 hostile mercenaries). The mutiny is always over bread.

\noindent\textbf{Clock}: \emph{Arrears [6]} – Ticks each time the party fails a negotiation or loses a fight. When full, the Captain abandons the contract. I have seen a captain leave in the middle of a siege; his men followed because the bread was stale.

\saikou{A contract is a leash. The best captains know which hand holds the other end. The worst captains forget and are hanged from their own standards.}

\subsection*{Linn Draugr‑Hunter}

\noindent\textbf{Description}: Northman who tracks draugr and other unquiet dead. Carries a silver‑edged axe, a bell, and a pouch of salt. They do not fear the dead. They fear dying without a witness.

\noindent\textbf{Stats}: Body 4, Wits 3, Spirit 4, Presence 2; Survival 3, Melee 3.\\
\textbf{Harm}: 5. \textbf{Fatigue}: Body 4.

\noindent\textbf{Traits}: \emph{Silver Net} – The Hunter may spend an action to throw a silver net at a spirit. The spirit must test Resolve (DV 4) or be immobilised for one exchange. The net is cold; it smells of the sea.

\noindent\textbf{Traits}: \emph{Bell‑Ring} – The Hunter rings a bell. For one scene, any spirit within Near has its Position worsened by one step. The bell’s note is the same note that tolled for the Hunter’s own grandfather.

\noindent\textbf{Clock}: \emph{Draugr’s Wake [6]} – Ticks each time the Hunter fails to put down a draugr. When full, a draugr lord rises and hunts the party. The lord remembers the Hunter’s name. It always remembers.

\dusana{The Linn sing when they hunt. The songs are not for the living. I have heard one. I will not write the tune.}

---

\section*{Rival Adventuring Parties}

\epigraph{“The road changes the walker. It changed them. It will change you.”}{—The Gray Wanderer, from a conversation I did not record}

Not every adversary is a monster. Some are mirrors. I have kept a list of the ones who walked past me in the bright stalls; I have kept another list of the ones who did not.

\subsection*{The Ashen Cord (Tier II–III)}

\noindent\textbf{Composition}: Vesper (leader, shadow mage, serves Ikasha), Tallow (berserker, serves Malachai), Ember (cantor, serves Palinode), Cinder (scout, no Patron). I met them once, in a rain‑soaked caravanserai. They were polite. They were watching.

\noindent\textbf{Goal}: Retrieve the same relic, secure a patron’s favor, or eliminate the PCs for a rival. Their goal is never stated; it is inferred from the empty seats they leave.

\noindent\textbf{Clock}: \emph{Cord’s Agenda [6]} – Ticks when the party ignores or fails to hinder them, when the GM spends 2+ SB on their off‑screen activity, or when a session ends without contact. When full, the Cord achieves a major milestone. I have seen their milestone. It was a closed door.

\noindent\textbf{Resource}: Hidden safehouse in Silkstrand (spend to hide from pursuers once). The safehouse has a red thread on the latch. Do not untie it.

\noindent\textbf{Relationship}: Unexpected rivals – will negotiate but may betray. They smile when they lie. It is a good smile.

\lerris{Vesper once offered me a cup of tea. I declined. The cup was warm. It should not have been.}

\subsection*{Custom Rival Generator}

Use the following d6 tables to generate a rival party in under 5 minutes. I have tested these tables; they produce plausible enemies, but they cannot produce the weight of a real grudge. For that, you must play.

\subsubsection*{Name (d6)}
\begin{enumerate}[label=\arabic*.]
    \item The Ashen Cord – echoes, smoke, unbreakable knots.
    \item The Brass Legion – gears, discipline, cold precision.
    \item The Thorns of Winter – ice, silence, patient cruelty.
    \item The Gilded Chain – wealth, obligation, gilded cages.
    \item The Salt Pact – oaths sworn in brine, pirates or pilgrims.
    \item The Unfaced – masks, secrets, stolen identities.
\end{enumerate}

\subsubsection*{Composition (d6)}
\begin{enumerate}[label=\arabic*.]
    \item Leader – shadow mage, war‑priest, strategist, or oath‑breaker.
    \item Muscle – berserker, knight, duelist, or beast‑handler.
    \item Mage/Weird – cantor, free caster, runekeeper, or witch.
    \item Scout/Thief – infiltrator, poisoner, sniper, or ghost.
    \item Face/Healer – negotiator, spy, medic, or confessor.
    \item Specialist – engineer, summoner, trapper, or lore‑master.
\end{enumerate}

\subsubsection*{Patron (d6 – roll once per party)}
\begin{enumerate}[label=\arabic*.]
    \item Ikasha (Shadow) – secrets, thresholds, omens.
    \item Malachai (Curse) – corrupted bargains, venomous gifts.
    \item The Traveler (Roads) – journeys, shortcuts, debts.
    \item The Ninth (Knowledge) – forbidden truths, fractured minds.
    \item Khemesh (Depths) – unstoppable pressure, drowning.
    \item Palinode (Performance) – rapture, encores, hungry fame.
\end{enumerate}

\subsubsection*{Goal (d6)}
\begin{enumerate}[label=\arabic*.]
    \item Retrieve the same artifact.
    \item Capture/kill a mutual contact.
    \item Secure a patron’s favor.
    \item Eliminate the PCs for a rival.
    \item Claim a contested site.
    \item Prevent a prophecy.
\end{enumerate}

\subsubsection*{Resource (d6)}
\begin{enumerate}[label=\arabic*.]
    \item Hidden safehouse – can hide from pursuers once.
    \item Bribed official – may delay a legal timer once.
    \item Spy network – gains +1 die to gather info.
    \item War chest – can hire temporary muscle.
    \item Patron’s omen – once per arc, avoid a Miss.
    \item Stolen map – shortcut to any location once.
\end{enumerate}

\gray{I have used this generator thrice. Twice, the rivals became allies. Once, they became graves. The graves were tidy.}

---

\section*{Terrestrial Patron Enforcers}

\epigraph{“A promise is a story we tell ourselves to make the world feel fair. The world is not fair.”}{—Dusana of the Raven Road, from a letter I found pressed in a ledger}

These adversaries collect for mortal masters – not cosmic patrons, but nobles, guilds, and syndicates with long memories and short tempers. They are the ones who knock when the payment is late.

\subsection*{Prefect’s Auditor (Ecktoria)}

\noindent\textbf{Description}: Imperial bureaucrat with legal authority to seize property. Carries a wax seal, a ledger, and a writ of seizure. The seal is always warm; the wax is always fresh.

\noindent\textbf{Stats}: Body 2, Wits 4, Spirit 3, Presence 3; Lore (law) 3, Sway 2.\\
\textbf{Harm}: 3. \textbf{Fatigue}: Body 2.

\noindent\textbf{Traits}: \emph{Legal Warfare} – The Auditor may spend 2 SB to impose a \emph{Seizure Timer } on a character’s Asset. When the timer fills, the Asset is confiscated. The notice is always polite; the bailiffs are not.

\noindent\textbf{Traits}: \emph{Immunity} – The Auditor cannot be directly attacked without triggering an immediate 4‑segment \emph{Arrest} timer. The law is a wall; the Auditor is the gate.

\noindent\textbf{Clock}: \emph{Seizure [4]} – Ticks each time the party fails a legal argument. When full, one Asset or minor magic item is taken. I have seen a man lose his grandmother’s ring to a clause he did not read.

\saikou{The law is a wall. A name is a door. The Auditor has many keys. The keys are not made of metal; they are made of precedent.}

\subsection*{Vermilion Corps Agent (Oshiira)}

\noindent\textbf{Description}: Internal investigator of the Oshiiran grain ledger. Wears red ink on their fingers and a smile like a trap. The ink is not for writing; it is for staining.

\noindent\textbf{Stats}: Body 3, Wits 4, Spirit 2, Presence 3; Investigation 3, Commerce 2.\\
\textbf{Harm}: 4. \textbf{Fatigue}: Body 3.

\noindent\textbf{Traits}: \emph{Grain Ledger} – The Agent may demand a tithe of grain (or its value in coin). Refusal advances the \emph{Embargo} timer. The tithe is never the same twice; it is calibrated to the weight of your silence.

\noindent\textbf{Traits}: \emph{Canal Lockdown} – Once per scene, the Agent may close a waterway, delaying travel by 2 segments. The lock’s key is a brass token shaped like a grain sheaf. I have not seen the token; I have only seen the delay.

\noindent\textbf{Clock}: \emph{Embargo [6]} – When full, the party cannot buy or sell grain in any Oshiiran port. The ports will still be open; the grain will be there; the price will be a memory.

\lerris{They once audited a ghost. The ghost paid in silence. The Agent filed the receipt under “settled”.}

\subsection*{Salt Prince’s Collector (Zakov)}

\noindent\textbf{Description}: Enforcer of the Salt Prince’s levy. Carries a brass hook, a ledger of debts, and a bottle of brine. The brine is not for drinking; it is for remembering.

\noindent\textbf{Stats}: Body 4, Wits 2, Spirit 3, Presence 2; Melee 2, Intimidation (Presence + Melee) 3.\\
\textbf{Harm}: 5. \textbf{Fatigue}: Body 4.

\noindent\textbf{Traits}: \emph{Debt Marker} – If the Collector has a physical token from a PC (hair, signature, a coin), they gain +1 die to track that PC across any distance. The token is kept in a pouch that pulses with a slow, wet rhythm.

\noindent\textbf{Traits}: \emph{Brine Curse} – Once per scene, the Collector may splash brine on a target. The target must test Resolve (DV 4) or gain the Condition \emph{Salt‑Tongue} (−1 die to all social rolls until they drink fresh water). The thirst is not for water; it is for the truth.

\noindent\textbf{Clock}: \emph{Levy Teeth [6]} – Ticks each time the party fails to pay tribute. When full, the Collector returns with 1d4 pirates and a writ of seizure. The pirates are always the same faces; the writ is always a different hand.

\gray{The Salt Prince’s crown floats. The Collector’s hook does not. I have seen the hook catch a debtor and pull him under. The water was not deep.}

---

\section*{Quick Adversary Cards (Cut out and use)}

The following six cards can be printed, cut, and kept behind the GM screen. They are not substitutes for attention; they are reminders that the world is full of people who owe and are owed.

\vspace{1cm}

\noindent
\fbox{\parbox{0.45\textwidth}{\centering \textbf{Bandit Captain} (TL 2)\\ Body 3, Melee 2, Harm 4\\ \textit{Rabble Rouser} – On a Hit, rally 1d4 bandits.\\ \textit{Cowardly} – Flees if outnumbered.}}
\hfill
\fbox{\parbox{0.45\textwidth}{\centering \textbf{Ghostly Anchor} (TL 3)\\ Spirit 4, Lore 2, Harm 3\\ \textit{Unfinished Business} – Cannot be harmed until its anchor is addressed.\\ \textit{Bargain} – May offer a deal.}}

\vspace{1cm}

\noindent
\fbox{\parbox{0.45\textwidth}{\centering \textbf{Slasher} (TL 3)\\ Body 4, Melee 3, Harm 5\\ \textit{Regeneration} – Ignores first 2 Harm per scene.\\ \textit{Hunger Timer [6]} – When full, frenzy.}}
\hfill
\fbox{\parbox{0.45\textwidth}{\centering \textbf{Oath‑Keeper} (TL –)\\ Cannot be fought.\\ \textit{Demand} – One story, memory, or confession.\\ \textit{Weakness} – Pay the original debt.}}

\vspace{1cm}

\noindent
\fbox{\parbox{0.45\textwidth}{\centering \textbf{Thorn Courtier} (TL 3)\\ Presence 4, Sway 3, Harm 3\\ \textit{Courteous Snare} – A gift accepted is a debt owed.\\ \textit{Iron Offense} – Brandishing iron gives the GM 2 SB.}}
\hfill
\fbox{\parbox{0.45\textwidth}{\centering \textbf{Salt Prince Enforcer} (TL 2)\\ Body 3, Melee 2, Harm 4\\ \textit{Debt Marker} – Knows where you sleep.\\ \textit{Brine Curse} – Test Resolve or gain Salt‑Tongue.}}

\vspace{1cm}

\lerris{I keep a card for myself. It reads: \textit{Lerris Fenwood – Debt: the ninth name. Payment: pending.} I do not recommend this practice.}

\gray{You’ll need more than cards. But it’s a start.}
\saikou{I keep a silver mirror in my pocket. It’s worth a hundred cards.}
\dusana{I keep a story. It’s worth a thousand.}
\lerris{I keep a ledger. The last page is for those who are not yet dead. I add a name each time I finish a chapter. This one is for Aqyl. He would have wanted it plain.}

